% ====================================================================
% §2.4 전이 중심의 온도 확장 — Einstein 진동 보정 (신설 절)
% ====================================================================
\section{전이 중심의 온도 확장 --- Einstein 진동 보정}\label{sec:einstein}

앞 절은 vibrational 엔트로피의 부분몰($x$-의존) 몫을 중심 표준값 $\Delta S^0_j$ 에 흡수시키되, 준양자 모드가
남기는 \emph{온도 편차}는 미결로 남겼다. 이 절은 그 편차를 파라미터 하나 --- 전이 $j$ 의 Einstein 온도
$\theta_{E,j}$ --- 로 명시적으로 회복해, 전이 중심 $U_j(T)$ 와 중심 엔트로피 $\Delta S^0_j$ 를 온도로
\emph{자기완결하게} 확장한다. 여기서 다루는 양은 $\theta_{E,j}$ 와 $x$ 를 고정하고 $T$ 만 바꾼 편차이지, 앞
절의 $\partial S_\vib/\partial x$ 가 아니다 --- 두 양의 혼동을 절 전체에서 피한다. 기호에 관해 먼저 못박아
둔다: Einstein 온도 $\theta_{E,j}$ 는 단위가 K 인 온도로, 앞 절들의 점유율 $\theta$(무차원)와 글자만 같은
전혀 다른 양이다 --- 아래첨자 $E$ 가 구분자이며, 본 절 이하에서 $\theta_E$ 표기는 항상 온도를 가리킨다.
이렇게 얻은 온도 확장은 이후
겹침 합성(\S\ref{sec:mixing})과 종합식(\S\ref{sec:synthesis})이 입력 $\{U_j(T),\Delta S^0_j\}$ 를 받을 때 그
자리에 그대로 들어간다.

\subsection{닫힌형 $S_\vib(T;\theta_{E,j})$ 와 두 극한}\label{ssec:einstein-closed}
준양자 모드 스펙트럼을 대표 진동수 하나로 접는다. 전이 $j$ 의 Einstein 온도를 $\theta_{E,j}\equiv\hbar
\omega_{E,j}/\kB$ 로 정의하면($\kB\theta_{E,j}\equiv\hbar\omega_{E,j}$), 무차원 비 $u_j\equiv\theta_{E,j}/T$ 로
점유가 $n=1/(e^{u_j}-1)$ 이다(비 $u_j$ 로 적어 본 장의 리튬 함량 $x$ 와 구분한다).\footnote{이 $u_j$ 는
앞 파트 \S\ref{sec:hys} spinodal 근의 $u_j=\sqrt{1-2RT/\Omega_j}$ 와도 동명 별개의 절-국소 기호다 --- 이 절
안에서 $u_j$ 는 오직 $\theta_{E,j}/T$ 를 뜻한다(부록 C 함정표 참조).} 대표값의 선택 --- 예컨대
LiCoO$_2$ 광학 포논 $50$--$80$~meV $\to\theta_{E,j}\approx580$--$930$~K --- 은 \emph{단위 환산 예시일 뿐}이며,
실제로는 대상 호스트의 실측 포논$\cdot$제일원리 포논 DOS \cite{jpcc2021} 로 교체하거나 다온도 곡률에서 직접
회귀한다[데이터-주도 1순위].

\textbf{(a) 출발 --- Einstein 단일 모드.} 식~\eqref{eq:Svib_mode} 유도의 모드 자유에너지
$f_k=\kB T\ln(1-e^{-\beta\hbar\omega_k})$ 에서 전 모드 합을 대표 진동수 하나로 접고, $S_\vib=-\partial
f_\vib/\partial T$ 를 1몰당으로 적용한다. \textbf{(b) 연산.} 점유 $n=1/(e^{u_j}-1)$ 을 모드 엔트로피
식~\eqref{eq:Svib_mode} 에 대입하면 $(1+n)\ln(1+n)-n\ln n$ 이 정리된다. \textbf{(c) 중간식.} 대수적으로
$(1+n)\ln(1+n)-n\ln n=-\ln(1-e^{-u_j})+u_j/(e^{u_j}-1)$ 이므로(단일모드 특수화이지 새 근사가 아니다),
\textbf{(d) 박스.}
\begin{equation}
\boxed{\;S_\vib(T;\theta_{E,j}) \;=\; R\Bigl[-\ln\bigl(1-e^{-u_j}\bigr)+\frac{u_j}{e^{u_j}-1}\Bigr],
\qquad u_j\equiv\frac{\theta_{E,j}}{T}\;.}
\label{eq:Svib-einstein}
\end{equation}
두 극한이 이 닫힌형을 검산한다. 고전극한 $\kB T\gg\hbar\omega_{E,j}$($u_j\to0$)에서
$-\ln(1-e^{-u_j})\to-\ln u_j+u_j/2$, $u_j/(e^{u_j}-1)\to1-u_j/2$ 가 상쇄해
\[
S_\vib\;\to\;R\bigl[1+\ln(T/\theta_{E,j})\bigr],
\]
곧 앞 절이 흡수한 ``$T$-무관 상수''는 닫힌형~\eqref{eq:Svib-einstein} 을 $T_\mathrm{ref}$ 에서 한 번 평가해
동결한 값 $S_\vib(T_\mathrm{ref};\theta_{E,j})$ 이며(그림~\ref{fig:svibid}(a) 캡션과 같은 서술), 고전 극한이
유효한 모드($\theta_{E,j}\ll T_\mathrm{ref}$)에서는 그 값이 위 보편식의 고온 코너
$R[1+\ln(T_\mathrm{ref}/\theta_{E,j})]$ 로 줄어든다 --- 단 본 절의 대표 준양자 작동점($\theta_{E,j}{\approx}700$~K,
$u_j{\approx}2.35$)에서는 닫힌형 $0.35R$ 대 고전 코너 $0.15R$ 로 갈리므로, 편차 정의~\eqref{eq:dSvib} 가 빼는
기준값은 언제나 닫힌형 쪽이다. 저온 $u_j\to\infty$ 에서는 $S_\vib\to0$(모든 모드가 바닥상태로 얼어붙어
흩어짐이 사라진다).

\textbf{적용 범위 --- 모드 수 셈과 부호.} 위 닫힌형은 전이당 \emph{알짜 한 모드}(진폭 $R$ 고정)가 반응에서
생성되는 셈법이다. 실제 반응 진동 엔트로피는 생성/반응 측 모드 집합의 \emph{차}이므로 모드 수가 균형이면
고전 극한에서 $T$-편차가 정확히 $0$ 으로 상쇄되는데(\S\ref{ssec:vib} 의 ``고전 극한 $=$ 상수 흡수'' 전제의
근거가 이 상쇄다), 단일모드 셈법의 편차~\eqref{eq:dSvib} 는 반대로 $u_j\to0$ 에서 $R\ln(T/T_\mathrm{ref})$ 로
가장 크고 기울기 $\partial\Delta S_\vib/\partial T=C_E(u_j)/T$($C_E$ 는 Einstein 비열, $0<C_E\le R$)가 항상
양이다. 곧 $\theta_{E,j}$ 지정은 알짜 연화형($\partial\Delta S_{\rxn,j}/\partial T>0$) 잔여를 갖는 준양자
전이에 한해 물리적 뜻이 있고, 경화형($\Delta C_p<0$) 잔여는 현행 단일모드 서식이 담지 않는다(부호 있는
진폭 또는 생성/반응 두-$\theta$ 차분형 일반화는 \emph{추가 후보}로 남긴다).

\subsection{기준온도 편차와 중심 이동 --- round-trip 정합}\label{ssec:einstein-roundtrip}
$\Delta S^0_j$ 는 $T_\mathrm{ref}$(기본 $298.15$~K)에서의 vib 기여를 이미 흡수했으므로, 새로 노출할 것은 그
기준 대비 편차뿐이다:
\begin{equation}
\Delta S_\vib(T) \;\equiv\; S_\vib(T;\theta_{E,j})-S_\vib(T_\mathrm{ref};\theta_{E,j}),
\qquad \Delta S_\vib(T_\mathrm{ref})=0.
\label{eq:dSvib}
\end{equation}
(이 $T$-편차 $\Delta S_\vib(T)$ 는 \S\ref{ssec:vib} 의 부분몰 $\partial S_\vib/\partial x$ 와 다른 양이다 ---
$\theta_{E,j}$ 와 $x$ 를 고정하고 $T$ 만 바꾼 것이다.) 이 편차를 전이 중심 표준값 자리에 가산해 유효 반응
엔트로피 $\Delta S_{\eff,j}(T)=\Delta S^0_j+\Delta S_\vib(T)$ 로 쓴다 --- 이 유효값이 이후 겹침 합성과
종합식의 $\Delta S^0_j$ 자리를 대신한다.

가산을 엔트로피 항 하나로만 하면 전위 중심과의 정합이 깨진다. 단순 곱 $T\Delta S_\vib/F$ 는 미분 시 여분
항을 남기기 때문이다. 대신 모드 자유에너지 편차에서 출발해 중심 이동을 닫는다.
\textbf{(a) 출발 --- 모드 자유에너지 편차.} 모드 Helmholtz 자유에너지 편차($T\!\to\!0$ 에서 $0$ 이 되는
열적 몫 --- 이 기준의 편차)를
\[
\Delta F_\vib(T)\;\equiv\; RT\ln\bigl(1-e^{-\theta_{E,j}/T}\bigr)
\]
로 두면(★ $F_\vib$ 는 몰당 Helmholtz 자유에너지 --- Faraday 상수 $F$ 와 구별), 엔트로피가 $S_\vib=-\partial
\Delta F_\vib/\partial T$ 로 자유에너지의 온도 미분으로 회복된다. \textbf{(b) 연산 --- round-trip 요구를
적분으로.} 구하려는 중심 이동은 $\partial\Delta U_\vib/\partial T=\Delta S_\vib(T)/F$ 를 만족해야 하므로,
편차 정의~\eqref{eq:dSvib} 를 $T_\mathrm{ref}$ 에서부터 적분한다:
\[
F\,\Delta U_\vib(T)=\int_{T_\mathrm{ref}}^{T}\Delta S_\vib(T')\,\dd T'
=\int_{T_\mathrm{ref}}^{T}S_\vib(T')\,\dd T'\;-\;S_\vib(T_\mathrm{ref})\,(T-T_\mathrm{ref}).
\]
\textbf{(c) 중간식 --- 적분을 자유에너지로 닫기.} 첫 적분은 $S_\vib=-\partial\Delta F_\vib/\partial T$
그대로 닫힌다:
\[
\int_{T_\mathrm{ref}}^{T}S_\vib(T')\,\dd T'
=-\bigl[\Delta F_\vib(T)-\Delta F_\vib(T_\mathrm{ref})\bigr].
\]
\textbf{(d) 박스.} 두 조각을 합치면 --- 정적분 하한이 $T_\mathrm{ref}$ 라 적분 상수
$\Delta U_\vib(T_\mathrm{ref})=0$ 은 자동으로 고정되고 ---
\begin{equation}
\boxed{\;
\Delta U_\vib(T)=-\frac1F\Bigl[\Delta F_\vib(T)-\Delta F_\vib(T_\mathrm{ref})+S_\vib(T_\mathrm{ref})\,(T-T_\mathrm{ref})\Bigr],
\qquad
\frac{\partial\,\Delta U_\vib}{\partial T}=\frac{\Delta S_\vib(T)}{F}\;}
\label{eq:dUvib}
\end{equation}
이다. 가산은 $U_j(T)\to U_j(T)+\Delta U_\vib(T)$ 와 $\Delta S^0_j\to\Delta S^0_j+\Delta S_\vib(T)$ 둘 다이며,
$\Delta U_\vib(T_\mathrm{ref})=0$ 이 자동이라 $T_\mathrm{ref}$ 에서 중심$\cdot$엔트로피가 보존된다. 미분 관계
$\partial\Delta U_\vib/\partial T=\Delta S_\vib(T)/F$ 가 곧 $\partial U_j/\partial T=\Delta S_{\eff,j}(T)/F$ 를 전위
항과 엔트로피 항에서 \emph{동시에} 만족하는 round-trip 정합이다 --- Bernardi 관계 $\partial U/\partial
T=\Delta S(T)/F$ 는 Kirchhoff $\Delta C_p$ 가 $U_\oc={-}\Delta G/F$ 미분에서 상쇄되어 $T$-의존 $\Delta S$ 에도
근사 없이 성립하므로, 중심과 발열을 \emph{같은} 자유에너지에서 짝지어 보정해야 하기 때문이다. 한편,
$\theta_{E,j}$ 를 지정하지 않은 전이는 두 가산이 모든 $T$ 에서 항등적으로 $0$ 이라 이 보정 이전과 완전히
동일한 값을 낸다(하위호환 --- 이 조건은 코드 개정의 회귀 기준으로 부록 D가 명세한다).

\subsection{수치 크기와 식별 --- 3온도점}\label{ssec:einstein-numid}
크기는 준양자 영역에서만 뜻이 있다. 대표값 $\theta_{E,j}=700$~K($\approx60$~meV)$\cdot$$T_\mathrm{ref}=298.15$~K
로 식~\eqref{eq:dUvib} 의 미분을 셈하면 $T=278.15,\,298.15,\,318.15,\,348.15$~K 에서 각각 $\partial\Delta
U_\vib/\partial T\approx-3.74,\,0,\,+3.70,\,+9.14~\mu$V/K 다 --- $T_\mathrm{ref}$ 를 지나며 부호가 바뀌고 온도
차가 커질수록 자란다. 식별의 근거는 함수형의 \emph{곡률}이다: $\Delta S_\vib(T;\theta_{E,j})$ 는 기울기와
곡률이 $\theta_{E,j}$ 하나에 묶인 비선형 1-모수족인 반면, electronic 잔여(식~\eqref{eq:Se-ch2}, $\propto T$)는
$T$ 에 선형이라 곡률이 $0$ 이다 --- 다온도 $dQ/dV$ 곡률 피팅이 두 함수형을 갈라 앞 절이 우려한 혼입을
해소한다. ($T_\mathrm{ref}$ 강제 영점 자체는 식별 근거가 아니다 --- \S\ref{ssec:decomp} 의 흡수 규약대로면
electronic 잔여도 $T_\mathrm{ref}$ 기준화되어 같은 영점을 공유하고, 기준화 여부와 무관하게 선형 성분은 자유
상수 $\Delta S^0_j$ 와 섞여 재정의될 뿐이다.) 단 분리에는 준양자 창 $\kB T\!\sim\!\hbar\omega$ 를 걸치는 $T$ 점
\textbf{셋 이상}이 필요하다 --- 2-온도 유한차분(국소 기울기만 봄)은 곡률과 선형을 합으로만 보아 축퇴하므로,
``섞일 수 있다''는 것은 저-온도점 진단이지 물리적 비식별성이 아니다. 닫힌형의 두 극한과 이 강제 영점·로그
곡률은 그림~\ref{fig:svibid} 가 수치 곡선으로 보인다. 진동 항의 온도 편차를 이렇게 닫았으니,
이제 세 분포를 겹침으로 조립할 준비가 되었다.

\begin{figure}[t]
\centering
\begin{tikzpicture}
% ============ (a) closed form S_vib/R vs T/theta_E ============
\begin{scope}[x=3.8cm,y=2.6cm]
\draw[-{Latex[length=1.6mm]}] (-0.02,0) -- (-0.02,1.62) node[above,font=\scriptsize] {$S_\vib/R$};
\draw[-{Latex[length=1.6mm]}] (-0.02,0) -- (1.72,0) node[below,font=\scriptsize] {$T/\theta_{E,j}$};
\foreach \xx in {0.5,1.0,1.5} {\draw (\xx,0.012) -- (\xx,-0.012) node[below,font=\scriptsize] {\xx};}
\foreach \yy in {0.5,1.0,1.5} {\draw (-0.008,\yy) -- (-0.032,\yy) node[left,font=\scriptsize] {\yy};}
% closed form (eq:Svib-einstein)
\draw[very thick,blue] plot[smooth] coordinates {(0.040,0.0000) (0.080,0.0001) (0.120,0.0022) (0.160,0.0140) (0.200,0.0407) (0.240,0.0812) (0.280,0.1318) (0.320,0.1885) (0.360,0.2484) (0.400,0.3092) (0.440,0.3698) (0.480,0.4293) (0.520,0.4872) (0.560,0.5433) (0.600,0.5974) (0.640,0.6496) (0.680,0.6998) (0.720,0.7482) (0.760,0.7947) (0.800,0.8395) (0.840,0.8827) (0.880,0.9243) (0.920,0.9644) (0.960,1.0032) (1.000,1.0407) (1.080,1.1119) (1.160,1.1788) (1.240,1.2418) (1.320,1.3012) (1.400,1.3575) (1.480,1.4108) (1.560,1.4616) (1.600,1.4861)};
% classical limit (dashed)
\draw[densely dashed] plot[smooth] coordinates {(0.380,0.0324) (0.420,0.1325) (0.460,0.2235) (0.500,0.3069) (0.540,0.3838) (0.580,0.4553) (0.620,0.5220) (0.660,0.5845) (0.700,0.6433) (0.740,0.6989) (0.780,0.7515) (0.820,0.8015) (0.860,0.8492) (0.900,0.8946) (0.940,0.9381) (0.980,0.9798) (1.060,1.0583) (1.140,1.1310) (1.220,1.1989) (1.300,1.2624) (1.380,1.3221) (1.460,1.3784) (1.540,1.4318) (1.600,1.4574)};
% labels (top-left is clear; dashed labeled near its low-T start)
\node[font=\tiny,anchor=west] at (0.06,1.50) {닫힌형 $S_\vib(T;\theta_{E,j})/R$ (식~\eqref{eq:Svib-einstein}, 실선)};
\node[font=\tiny,anchor=west] at (0.43,0.045) {고전 극한 $1+\ln(T/\theta_E)$ (파선)};
% freeze-out annotation
\node[font=\tiny,anchor=west,align=left] at (0.05,0.30) {저온 동결\\$S_\vib\to0$};
\draw[-{Latex[length=1.2mm]},black!60] (0.10,0.24) -- (0.115,0.02);
% operating point theta_E=700K, T=298K
\fill (0.426,0.3486) circle (1.3pt);
\node[font=\tiny,anchor=west,align=left] at (0.52,0.30) {$\theta_E{=}700$ K, $T{=}298$ K\\(준양자 작동점, $S_\vib/R{=}0.35$)};
\draw[-{Latex[length=1.2mm]},black!60] (0.51,0.33) -- (0.443,0.345);
\node[font=\scriptsize] at (0.82,-0.22) {(a)};
\end{scope}
% ============ (b) forced zero of the center-shift slope ============
\begin{scope}[xshift=8.1cm,x=0.060cm,y=0.165cm,shift={(-258.15,0)}]
\draw[-{Latex[length=1.6mm]}] (256,-10.5) -- (256,14.6)
  node[above,font=\scriptsize] {$\partial\Delta U_\vib/\partial T=\Delta S_\vib(T)/F$ [$\mu$V/K]};
\draw[-{Latex[length=1.6mm]}] (256,-10.5) -- (362,-10.5) node[anchor=north west,xshift=-2mm,font=\scriptsize] {$T$ [K]};
\foreach \xx in {260,280,300,320,340} {\draw (\xx,-10.2) -- (\xx,-10.8) node[below,font=\scriptsize] {\xx};}
\foreach \yy in {-8,-4,0,4,8,12} {\draw (256.4,\yy) -- (255.6,\yy) node[left,font=\scriptsize] {\yy};}
\draw[densely dashed,black!45] (256,0) -- (360,0);
\draw[densely dotted,black!55] (298.15,-10.5) -- (298.15,12.2);
\node[font=\tiny,black!55,anchor=west] at (300.5,-6.0) {$T_\mathrm{ref}$ 강제 영점};
\draw[-{Latex[length=1.1mm]},black!55] (302,-4.9) -- (298.6,-0.5);
% theta_E = 500 K (dashed)
\draw[thick,densely dashed] plot[smooth] coordinates {(258.15,-9.530) (263.15,-8.303) (268.15,-7.086) (273.15,-5.879) (278.15,-4.682) (283.15,-3.496) (288.15,-2.320) (293.15,-1.155) (298.15,0.000) (303.15,1.144) (308.15,2.277) (313.15,3.400) (318.15,4.512) (323.15,5.614) (328.15,6.705) (333.15,7.786) (338.15,8.856) (343.15,9.916) (348.15,10.967) (353.15,12.007) (358.15,13.037)};
% theta_E = 700 K (solid blue)
\draw[very thick,blue] plot[smooth] coordinates {(258.15,-7.487) (263.15,-6.551) (268.15,-5.614) (273.15,-4.676) (278.15,-3.738) (283.15,-2.802) (288.15,-1.866) (293.15,-0.932) (298.15,0.000) (303.15,0.930) (308.15,1.856) (313.15,2.780) (318.15,3.700) (323.15,4.617) (328.15,5.530) (333.15,6.438) (338.15,7.343) (343.15,8.243) (348.15,9.138) (353.15,10.029) (358.15,10.915)};
% theta_E = 900 K (dotted)
\draw[thick,densely dotted] plot[smooth] coordinates {(258.15,-5.516) (263.15,-4.851) (268.15,-4.177) (273.15,-3.496) (278.15,-2.808) (283.15,-2.114) (288.15,-1.414) (293.15,-0.709) (298.15,0.000) (303.15,0.713) (308.15,1.430) (313.15,2.149) (318.15,2.872) (323.15,3.596) (328.15,4.322) (333.15,5.049) (338.15,5.777) (343.15,6.506) (348.15,7.235) (353.15,7.965) (358.15,8.694)};
% text anchors on the 700 K curve (values from §2.4.3) — white-backed tags
\fill (278.15,-3.738) circle (1.2pt);
\node[font=\tiny,fill=white,inner sep=0.8pt] at (281.5,-4.9) {$-3.74$};
\fill (298.15,0.000) circle (1.2pt);
\fill (318.15,3.700) circle (1.2pt);
\node[font=\tiny,fill=white,inner sep=0.8pt] at (321.5,2.55) {$+3.70$};
\fill (348.15,9.138) circle (1.2pt);
\node[font=\tiny,fill=white,inner sep=0.8pt] at (351.3,8.0) {$+9.14$};
% legend (top-left clear)
\draw[thick,densely dashed] (259,13.4) -- (266,13.4);
\node[font=\tiny,anchor=west] at (267,13.4) {$\theta_E{=}500$ K};
\draw[very thick,blue] (259,11.8) -- (266,11.8);
\node[font=\tiny,anchor=west] at (267,11.8) {$700$ K};
\draw[thick,densely dotted] (259,10.2) -- (266,10.2);
\node[font=\tiny,anchor=west] at (267,10.2) {$900$ K};
\node[font=\scriptsize] at (308,-13.6) {(b)};
\end{scope}
\end{tikzpicture}
\caption{Einstein 진동 보정의 두 얼굴 --- 좌표는 식~\eqref{eq:Svib-einstein}$\cdot$\eqref{eq:dUvib}
그대로의 수치 평가다. (a) 닫힌형 $S_\vib(T;\theta_{E,j})/R$: 저온($u_j\to\infty$)에서 $0$ 으로
동결되고 고전 극한(파선) $1+\ln(T/\theta_E)$ 에 $T/\theta_E\gtrsim1$ 에서 붙는다 --- 앞 절이 중심
표준값에 흡수한 ``$T$-무관 상수''는 이 곡선을 $T_\mathrm{ref}$ 에서 한 번 평가해 동결한 값이다.
● 는 $\theta_E{=}700$ K$\cdot$$298$ K 의 준양자 작동점. (b) 중심 이동의 온도 기울기
$\partial\Delta U_\vib/\partial T=\Delta S_\vib(T)/F$: $\theta_{E}{=}500/700/900$ K 모두
$T_\mathrm{ref}$ 의 \emph{강제 영점}을 지나 로그형 곡률로 자라며, ● 는 본문 \S\ref{ssec:einstein-numid}
의 수치($\theta_E{=}700$ K 에서 $-3.74/0/+3.70/+9.14\ \mu$V/K)다. electronic 항($\propto T$)은 곡률이 $0$
이라 두 함수형이 다온도 곡률 피팅에서 갈리고, 분리에는 준양자 영역에 놓인 온도점 셋 이상이 필요하다.}
\label{fig:svibid}
\end{figure}
% 자산(v1.0.21 신규): [V21-R-05 Einstein 닫힌형·강제 영점 2패널 fig:svibid — FIGS_PICK N-4(베이스 FF2-5)]

% 자산: [C-49] θ_E,j≡ℏω_E,j/k_B·LCO 50-80meV→580-930K 환산 예시일 뿐(데이터 1순위) ·
%   [C-50] f_k 대표진동수 접기·u_j≡θ_E,j/T(x와 구분) · [C-51] ★S_vib(T;θ_E,j) (eq:Svib-einstein) ·
%   [C-52] 고온극한 R[1+ln(T/θ_E)]=동결 상수 · [C-53] 저온극한 S_vib→0 ·
%   [C-54] ΔS_vib(T)≡S_vib(T)−S_vib(T_ref), T_ref서 0 (eq:dSvib) · [C-55] ΔS_vib(T)≠∂S_vib/∂x 구분 ·
%   [C-56] ΔS_eff,j(T)=ΔS⁰_j+ΔS_vib(T) 가산(입력 자리) · [C-57] 단순곱 여분항→F_vib(Helmholtz≠Faraday F) ·
%   [C-58] ★ΔU_vib(T)·∂ΔU_vib/∂T=ΔS_vib/F (eq:dUvib) · [C-59] 가산 U_j·ΔS⁰_j 둘 다·T_ref서 0 자동 ·
%   [C-60] round-trip 근거(Kirchhoff ΔC_p 상쇄) · [C-61] 하위호환 θ_E 미지정 두 가산 0(회귀 기준→부록B) ·
%   [C-62] 수치 −3.74/0/+3.70/+9.14 µV/K(θ_E=700K·T_ref=298.15·4온도점) · [C-63] ★식별 elec(∝T)vs vib(영점+로그)·3온도점 필요
